\documentclass[11pt]{article}
\usepackage[a4paper,margin=1in]{geometry}
\usepackage{fontspec}
\usepackage[boldfont=false]{xeCJK}
\setCJKmainfont{FandolSong-Regular.otf}
\usepackage{amsmath}
\usepackage{amssymb}
\usepackage{array}
\usepackage{booktabs}
\usepackage{graphicx}
\usepackage{adjustbox}
\usepackage{enumitem}
\setlist[itemize]{topsep=2pt,partopsep=0pt,parsep=0pt,itemsep=2pt,leftmargin=1.4em}
\setlist[enumerate]{topsep=2pt,partopsep=0pt,parsep=0pt,itemsep=2pt,leftmargin=1.6em}
\usepackage{parskip}
\usepackage{fvextra}
\fvset{breaklines=true,breakanywhere=true,fontsize=\small}
\setlength{\parindent}{0pt}

\begin{document}

\begin{center}
  {\LARGE\bfseries Organic chemistry}\\[0.35em]
  {\large Igcse Chemistry}
\end{center}
\vspace{0.6em}

Organic chemistry is the chemistry of carbon compounds. Carbon is special because it can join to other carbon atoms to make long chains and rings.

\section*{Formulae and key words}

There are several ways to write an organic molecule:

\begin{itemize}
\item The \textbf{molecular formula} 分子式 shows how many of each atom there are, for example $\text{C}_2\text{H}_6$.

\item The \textbf{displayed formula} 结构式 shows every atom and every bond drawn out in full.

\item The \textbf{structural formula} 结构简式 shows how the atoms are arranged without drawing every bond, for example $\text{CH}_3\text{CH}_2\text{OH}$.

\item The \textbf{general formula} 通式 works for a whole family, for example $\text{C}_n\text{H}_{2n+2}$ for alkanes.

\end{itemize}

A \textbf{homologous series} 同系物 is a family of compounds with the same \textbf{functional group} 官能团 — the atom or group of atoms that gives the family its chemical properties. Members of a series have the same general formula, differ by a $\text{CH}_2$ unit each step, and have similar chemical properties with a gradual change in physical properties.

\begin{center}
\adjustbox{max width=\linewidth}{%
\begin{tabular}{lll}
\toprule
\textbf{Series} & \textbf{Functional group} & \textbf{General formula} \\
\midrule
alkanes & C–C single bonds only & $\text{C}_n\text{H}_{2n+2}$ \\
alkenes & C=C double bond & $\text{C}_n\text{H}_{2n}$ \\
alcohols & –OH & $\text{C}_n\text{H}_{2n+1}\text{OH}$ \\
carboxylic acids & –COOH & $\text{C}_n\text{H}_{2n+1}\text{COOH}$ \\
\bottomrule
\end{tabular}}
\end{center}

A \textbf{saturated} 饱和 compound has only single carbon–carbon bonds. An \textbf{unsaturated} 不饱和 compound has one or more carbon–carbon bonds that are not single (such as a C=C double bond).

\textbf{Structural isomers} 同分异构体 are compounds with the same molecular formula but different structural formulae. For example, $\text{C}_4\text{H}_{10}$ can be a straight chain or a branched chain.

\section*{Naming organic compounds}

The end of the name tells you the family:

\begin{center}
\adjustbox{max width=\linewidth}{%
\begin{tabular}{lll}
\toprule
\textbf{Ending} & \textbf{Family} & \textbf{Example} \\
\midrule
-ane & alkane & methane, ethane \\
-ene & alkene & ethene \\
-ol & alcohol & ethanol \\
-oic acid & carboxylic acid & ethanoic acid \\
\bottomrule
\end{tabular}}
\end{center}

The start of the name tells you the number of carbon atoms: \emph{meth-} = 1, \emph{eth-} = 2, \emph{prop-} = 3, \emph{but-} = 4. For longer alkenes and alcohols, a number shows where the functional group is, for example but-1-ene and but-2-ene, or propan-1-ol and propan-2-ol.

\section*{Fuels}

The three \textbf{fossil fuels} 化石燃料 are \textbf{coal} 煤, \textbf{natural gas} 天然气 and \textbf{petroleum} 石油 (crude oil). \textbf{Methane} 甲烷 is the main part of natural gas.

A \textbf{hydrocarbon} 碳氢化合物 is a compound made of hydrogen and carbon only. Petroleum is a mixture of many different hydrocarbons.

\subsection*{Fractional distillation}

Petroleum is separated into useful \textbf{fractions} 馏分 by \textbf{fractional distillation} 分馏. The mixture is heated, and the different hydrocarbons turn to gas and rise up a tall \textbf{fractionating column} 分馏塔. The column is hot at the bottom and cool at the top, so each fraction turns back to liquid at a different height.

Going from the bottom to the top of the column, the fractions have:

\begin{itemize}
\item shorter \textbf{chain length} 链长 (smaller molecules);

\item higher \textbf{volatility} 挥发性 (they turn to gas more easily);

\item lower boiling points;

\item lower \textbf{viscosity} 黏度 (they flow more easily).

\end{itemize}

\begin{center}
\adjustbox{max width=\linewidth}{%
\begin{tabular}{ll}
\toprule
\textbf{Fraction} & \textbf{Use} \\
\midrule
refinery gas & gas for heating and cooking \\
gasoline / petrol & fuel for cars \\
naphtha & raw material for making chemicals \\
kerosene / paraffin & jet fuel \\
diesel oil & fuel for diesel engines \\
fuel oil & fuel for ships and home heating \\
lubricating oil & lubricants, waxes and polishes \\
\textbf{bitumen} 沥青 & making roads \\
\bottomrule
\end{tabular}}
\end{center}

\section*{Alkanes}

\textbf{Alkanes} 烷烃 have only single covalent bonds, so they are saturated hydrocarbons. They are generally unreactive. Their two important reactions are:

\begin{itemize}
\item \textbf{Combustion} 燃烧: they burn in plenty of oxygen to give carbon dioxide and water.

\item \textbf{Substitution} with chlorine.

\end{itemize}

In a \textbf{substitution reaction} 取代反应, one atom (or group of atoms) is replaced by another. Alkanes react with chlorine only in \textbf{ultraviolet light} 紫外线 — this is a \textbf{photochemical reaction} 光化学反应, where the light provides the \textbf{activation energy} 活化能. For example:

\[
\text{CH}_4 + \text{Cl}_2 \rightarrow \text{CH}_3\text{Cl} + \text{HCl}
\]

\section*{Alkenes}

\textbf{Alkenes} 烯烃 have a carbon–carbon double bond (C=C), so they are unsaturated hydrocarbons.

\subsection*{Cracking}

Large alkane molecules are not very useful. \textbf{Cracking} 裂化 breaks them into smaller, more useful molecules — smaller alkanes and alkenes — using a high temperature and a \textbf{catalyst} 催化剂. Cracking also makes \textbf{hydrogen} 氢气 and provides alkenes for making plastics.

\subsection*{Reactions of alkenes}

The C=C double bond makes alkenes reactive. They take part in \textbf{addition reactions} 加成反应, where two molecules join to form a single product.

\begin{itemize}
\item \textbf{Test for unsaturation}: shake the compound with \textbf{bromine} 溴 water (which is orange). An alkene turns the bromine water colourless; an alkane does not change it.

\item With hydrogen and a \textbf{nickel} 镍 catalyst, an alkene becomes an alkane.

\item With \textbf{steam} 蒸汽 and an acid catalyst, an alkene becomes an alcohol.

\end{itemize}

\section*{Alcohols}

\textbf{Alcohols} 醇 contain the –OH functional group. The most important one is \textbf{ethanol} 乙醇. There are two ways to make ethanol.

\begin{center}
\adjustbox{max width=\linewidth}{%
\begin{tabular}{lll}
\toprule
\textbf{Method} & \textbf{Conditions} & \textbf{Notes} \\
\midrule
\textbf{fermentation} 发酵 of glucose & yeast, 25–35 °C, no oxygen & uses renewable sugar, but slow and gives impure ethanol \\
addition of steam to \textbf{ethene} 乙烯 & 300 °C, 60 atm, acid catalyst & fast and pure, but uses petroleum (non-renewable) \\
\bottomrule
\end{tabular}}
\end{center}

In \textbf{fermentation}, \textbf{yeast} 酵母 turns \textbf{glucose} 葡萄糖 into ethanol and carbon dioxide.

Ethanol burns well (combustion), so it is used as a \textbf{fuel}. It also dissolves many substances, so it is used as a \textbf{solvent} 溶剂.

\section*{Carboxylic acids}

\textbf{Carboxylic acids} 羧酸 contain the –COOH functional group. \textbf{Ethanoic acid} 乙酸 is the one to know. Like other acids, it reacts with:

\begin{itemize}
\item \textbf{metals} 金属 → a salt + hydrogen;

\item \textbf{bases} 碱 → a \textbf{salt} 盐 + water;

\item \textbf{carbonates} 碳酸盐 → a salt + water + carbon dioxide.

\end{itemize}

The salts formed are called ethanoates.

Ethanoic acid can be made by the \textbf{oxidation} 氧化 of ethanol, either using acidified \textbf{potassium manganate(VII)} 高锰酸钾, or by bacteria during the making of \textbf{vinegar} 醋.

When a carboxylic acid reacts with an alcohol (using an acid catalyst), it forms an \textbf{ester} 酯.

\section*{Polymers}

A \textbf{polymer} 聚合物 is a very large molecule built from many small molecules called \textbf{monomers} 单体.

\subsection*{Addition polymerisation}

In \textbf{addition polymerisation} 加聚, many unsaturated monomers join together with no other product. For example, many ethene monomers join to make poly(ethene). The small part that repeats along the chain is the \textbf{repeat unit} 重复单元.

\subsection*{Condensation polymerisation}

In \textbf{condensation polymerisation} 缩聚, monomers join and a small molecule (usually water) is lost each time a bond forms. This makes two important types:

\begin{itemize}
\item \textbf{Polyamides} 聚酰胺 are made from a dicarboxylic acid and a diamine. \textbf{Nylon} 尼龙 is a polyamide.

\item \textbf{Polyesters} 聚酯 are made from a dicarboxylic acid and a diol. PET is a polyester, and it can be broken back into its monomers and re-made.

\end{itemize}

\subsection*{Plastics and the environment}

\textbf{Plastics} 塑料 are made from polymers. Because many plastics do not break down, getting rid of them is a problem:

\begin{itemize}
\item they build up in \textbf{landfill} 填埋 sites;

\item they collect in the oceans and harm sea life;

\item burning them can make toxic gases.

\end{itemize}

\subsection*{Proteins}

\textbf{Proteins} 蛋白质 are natural polyamides. They are made from \textbf{amino acid} 氨基酸 monomers joined in long chains.


\end{document}
